%!TEX root = ../main.tex
\chapter{Zdecentralizowana infrastruktura klucza publicznego}
\label{ch:dpki}

\section{Koncepcja DPKI}
\label{sec:koncepcja-dpki}

Idea infrastruktury klucza publicznego pozbawionej centralnego pośrednika jest starsza niż blockchain. Już sieć zaufania PGP \cite{zimmermann1995pgp} proponowała model, w którym wiarygodność klucza wynika nie z podpisu urzędu, lecz z grafu wzajemnych poświadczeń użytkowników. Podejście to nigdy jednak nie znalazło zastosowania w skali Web PKI - z trzech powodów, które warto odnotować, bo wyznaczają one wymagania wobec każdego następcy: brak globalnie spójnego widoku (różni weryfikujący mogą dojść do różnych wniosków o tym samym kluczu), brak systemowego mechanizmu odnajdywania kluczy oraz praktycznie niedziałające unieważnianie.

Termin \emph{Decentralized Public Key Infrastructure} w obecnym rozumieniu ukształtował się wraz z dojrzewaniem technologii blockchain \cite{allen2015dpki}. Sednem koncepcji jest zastąpienie poświadczenia zaufanej trzeciej strony zapisem w rozproszonym rejestrze: wiązanie identyfikatora (w kontekście tej pracy - nazwy domeny) z kluczem publicznym jest tworzone, aktualizowane i odwoływane transakcjami podpisywanymi kluczem właściciela, a o przyjęciu transakcji decyduje protokół konsensusu egzekwujący jawne, jednakowe dla wszystkich reguły. Rejestr pełni przy tym trzy role naraz, które w klasycznej PKI są rozproszone między CA, repozytoria certyfikatów i infrastrukturę CRL/OCSP:
\begin{itemize}
    \item \textbf{źródła prawdy} - aktualny stan wiązania (w tym jego unieważnienie) jest częścią globalnego, replikowanego stanu, a nie odpowiedzią pojedynczego serwera;
    \item \textbf{dziennika zdarzeń} - pełna historia operacji na identyfikatorze jest trwała i publicznie audytowalna, co zapewnia transparentność silniejszą niż dobrowolne logi CT;
    \item \textbf{punktu egzekwowania polityki} - reguły takie jak ,,tylko posiadacz aktualnego klucza może rotować lub unieważnić wpis'' są wymuszane przez konsensus, a nie przez procedury operacyjne zaufanego podmiotu.
\end{itemize}

Kluczową konsekwencją jest odwrócenie kierunku kontroli: w modelu hierarchicznym to urząd dysponuje władzą nad wiązaniem (wystawia je i unieważnia), w DPKI władza ta należy wyłącznie do posiadacza klucza prywatnego identyfikatora. Eliminuje to klasę ataków przez kompromitację pośrednika, ale zarazem przenosi pełną odpowiedzialność za klucz na właściciela - utrata klucza bez wcześniejszej rotacji oznacza utratę kontroli nad identyfikatorem, a żaden podmiot nie ma uprawnień, by ją przywrócić. Otwartym problemem każdej DPKI pozostaje też inicjalne związanie identyfikatora z podmiotem: rejestr gwarantuje, że wpisu nie da się podrobić ani zmodyfikować, ale to protokół rejestracji musi zagwarantować, że trafia do niego wpis właściwego podmiotu. Dla identyfikatorów zakorzenionych w istniejących systemach (jak nazwy domen w DNS) naturalnym rozwiązaniem jest weryfikacja przez wyzwanie, analogiczna do ACME; sposób jej pogodzenia z deterministycznym konsensusem omówiono w rozdziale~\ref{ch:projekt}.

\section{Self-Sovereign Identity i DID}
\label{sec:ssi-did}

Równolegle do prac nad zdecentralizowaną PKI dla domen rozwinął się nurt tożsamości suwerennej (\emph{Self-Sovereign Identity}, SSI), przenoszący tę samą zasadę - identyfikator kontrolowany wyłącznie kluczem właściciela - na tożsamości osób i organizacji. Jego standaryzacyjnym owocem są zdecentralizowane identyfikatory DID \cite{w3c-did-core}: identyfikatory URI, które rozwiązują się do dokumentu DID zawierającego m.in.\ klucze publiczne i punkty usługowe podmiotu, przy czym sposób rozwiązywania definiuje tzw.\ metoda DID - w wielu metodach rolę weryfikowalnego rejestru danych pełni właśnie blockchain. Nad warstwą identyfikatorów nadbudowywane są poświadczenia weryfikowalne (\emph{Verifiable Credentials}), pozwalające przenosić podpisane atrybuty tożsamości między stronami bez udziału wystawcy.

Z perspektywy niniejszej pracy nurt SSI jest istotny z dwóch powodów. Po pierwsze, potwierdza on architektoniczny wzorzec wspólny dla wszystkich systemów tej klasy: rozproszony rejestr jako korzeń zaufania, operacje na tożsamości jako podpisane transakcje, weryfikacja po stronie klienta na podstawie dowodów kryptograficznych. Po drugie, wyznacza granicę zakresu pracy: SSI adresuje tożsamość \emph{podmiotów} i kładzie nacisk na prywatność, selektywne ujawnianie atrybutów i relacje między wystawcami poświadczeń, podczas gdy Web PKI dotyczy \emph{publicznych z natury} wiązań domena--klucz, gdzie zagadnienia prywatności ustępują miejsca spójności globalnego widoku i natychmiastowości unieważnień. Projektowany system jest więc pokrewny metodom DID na poziomie architektury, lecz celowo zawężony do problemu, w którym słabości klasycznej PKI są najlepiej udokumentowane.

\section{Przegląd istniejących rozwiązań}
\label{sec:przeglad-rozwiazan}

\subsection{Namecoin}
\label{subsec:namecoin}

Namecoin, uruchomiony w 2011 roku jako pierwszy fork Bitcoina, jest historycznie pierwszą realizacją zdecentralizowanego rejestru nazw. System utrzymuje przestrzeń nazw klucz--wartość (najbardziej znana jest domena najwyższego poziomu \texttt{.bit}) z rejestracją w modelu ,,kto pierwszy, ten lepszy'', okresowym odnawianiem wpisów oraz aktualizacją wyłącznie przez posiadacza klucza. Konsensus dziedziczony jest z Bitcoina poprzez kopanie scalone (\emph{merged mining}): górnicy Bitcoina mogą bez dodatkowego kosztu zabezpieczać równocześnie łańcuch Namecoina.

Empiryczna analiza systemu \cite{kalodner2015namecoin} obnażyła jednak zasadnicze problemy. Przestrzeń nazw została zdominowana przez spekulacyjne rezerwacje - spośród około 120 tysięcy zarejestrowanych nazw autorzy zidentyfikowali zaledwie 28 niosących nietrywialną treść; model ,,kto pierwszy, ten lepszy'' bez weryfikacji jakiegokolwiek związku rejestrującego z nazwą okazał się zaproszeniem do masowego przejmowania nazw wartościowych. Jeszcze poważniejszy okazał się skutek uboczny kopania scalonego: koncentracja mocy obliczeniowej doprowadziła do sytuacji, w której pojedyncza kopalnia przekraczała połowę mocy sieci, unieważniając podstawowe założenie bezpieczeństwa PoW. Wreszcie probabilistyczna finalność i pełnowęzłowy model weryfikacji czyniły system nieprzyjaznym dla lekkich klientów. Dla niniejszej pracy Namecoin jest cenną lekcją negatywną w trzech punktach: rejestracja musi być powiązana z weryfikacją kontroli nad identyfikatorem, bezpieczeństwo konsensusu nie może być efektem ubocznym innego systemu, a weryfikacja po stronie klienta musi być możliwa bez utrzymywania pełnej kopii łańcucha.

\subsection{Blockstack/Stacks}
\label{subsec:blockstack}

Blockstack \cite{ali2016blockstack} stanowi drugą generację podejścia. Po doświadczeniach z koncentracją mocy w Namecoinie (system pierwotnie działał właśnie na nim) autorzy przenieśli rejestr nazw na łańcuch Bitcoina, wprowadzając warstwę \emph{virtualchain}: logika systemu nazw (BNS) jest osadzana w zwykłych transakcjach Bitcoina i interpretowana przez węzły Blockstacka, dzięki czemu rejestr dziedziczy bezpieczeństwo najsilniejszego łańcucha PoW. Drugim wkładem jest konsekwentne rozdzielenie płaszczyzny kontroli od płaszczyzny danych: w łańcuchu zapisywane są wyłącznie skróty i wskaźniki, a właściwe dane (odpowiedniki stref DNS, profile) przechowywane są poza łańcuchem w systemach magazynowych, których integralność weryfikowana jest względem skrótów z rejestru. W późniejszej ewolucji projekt przekształcił się w samodzielny łańcuch Stacks z własnym mechanizmem konsensusu zakotwiczonym w Bitcoinie.

Ograniczenia tego podejścia są lustrzanym odbiciem jego zalet. Dziedzicząc bezpieczeństwo Bitcoina, system dziedziczy również jego probabilistyczną finalność oraz przepustowość i opóźnienia - operacje na nazwach konkurują o miejsce w blokach z transakcjami finansowymi, a potwierdzenie rejestracji wymaga oczekiwania liczonego w dziesiątkach minut. Rejestracja pozostaje przy tym, jak w Namecoinie, pozbawiona weryfikacji związku rejestrującego z nazwą, co czyni system odrębną przestrzenią nazw, a nie warstwą zaufania dla istniejącego DNS - adopcja wymaga zmiany sposobu adresowania, nie tylko sposobu weryfikacji kluczy.

\subsection{CertCoin i prace pokrewne}
\label{subsec:certcoin}

CertCoin \cite{fromknecht2014certcoin} jest pierwszą propozycją wprost ukierunkowaną na zastąpienie PKI: na bazie Namecoina definiuje pełny cykl życia wiązania domena--klucz (rejestracja, aktualizacja, weryfikacja, unieważnienie), z parą kluczy online/offline ułatwiającą odzyskiwanie po kompromitacji oraz - co najistotniejsze - z akumulatorami kryptograficznymi i strukturami uwierzytelnionymi umożliwiającymi zwięzłe dowody dla lekkich klientów. CertCoin pozostał jednak konstrukcją koncepcyjną: dziedziczy wszystkie omówione słabości podłoża (PoW, probabilistyczna finalność, rejestracja bez weryfikacji własności), a jego ewaluacja nie wykroczyła poza prototyp komponentów.

Z licznych prac pokrewnych warto wyróżnić dwie, reprezentujące odmienne filozofie. IKP \cite{matsumoto2017ikp} nie zastępuje CA, lecz nadbudowuje nad nimi warstwę odpowiedzialności: właściciele domen publikują w łańcuchu (Ethereum) polityki certyfikacyjne, a wystawienie certyfikatu z nimi sprzecznego uruchamia automatyczne kary finansowe dla CA i nagrody dla zgłaszających. Podejście to wzmacnia ekonomiczną dyscyplinę istniejącego ekosystemu, ale z założenia zachowuje jego model zaufania. CONIKS \cite{melara2015coniks} z kolei, choć dotyczy kluczy użytkowników końcowych, a nie domen, dostarczył konstrukcji o fundamentalnym znaczeniu dla tej pracy: uwierzytelnionego słownika opartego na drzewie Merkle z dowodami zarówno obecności, jak i \emph{nieobecności} klucza, publikowanego okresowo przez dostawcę i audytowalnego przez klientów. Brakującym elementem CONIKS jest rozproszone uzgadnianie publikowanego korzenia - pojedynczy dostawca może próbować rozszczepiać widok (\emph{equivocation}), a wykrywalność tego zależy od zewnętrznych mechanizmów; późniejsze prace proponowały zakotwiczenie korzeni CONIKS w łańcuchu Ethereum w celu domknięcia tej luki \cite{bonneau2016ethiks}.

\medskip

Zestawienie powyższych systemów ujawnia spójny obraz luki badawczej. Żadne z omówionych rozwiązań nie łączy jednocześnie: (1) deterministycznej finalności konsensusu, bez której status certyfikatu nigdy nie jest ostateczny, (2) weryfikacji kontroli nad istniejącym identyfikatorem DNS w momencie rejestracji, bez której rejestr chroni jedynie własną, odrębną przestrzeń nazw, (3) dowodliwych odpowiedzi o stan - włącznie z dowodem nieobecności i unieważnienia - weryfikowalnych przez lekkiego klienta bez zaufania do odpytywanego węzła, oraz (4) działającej implementacji poddanej systematycznej analizie bezpieczeństwa i wydajności. Konstrukcja systemu spełniającego łącznie te cztery wymagania - z konsensusem BFT (sekcja~\ref{subsec:bft-tendermint}) w miejsce PoW, weryfikacją DNS-01 wzorowaną na ACME w miejsce rejestracji ,,kto pierwszy, ten lepszy'' oraz rzadkim drzewem Merkle (sekcja~\ref{sec:merkle}) w roli stanu aplikacji - jest przedmiotem pozostałej części pracy.
